% Options for packages loaded elsewhere
\PassOptionsToPackage{unicode}{hyperref}
\PassOptionsToPackage{hyphens}{url}
\documentclass[
]{article}
\usepackage{xcolor}
\usepackage[margin=1in]{geometry}
\usepackage{amsmath,amssymb}
\setcounter{secnumdepth}{-\maxdimen} % remove section numbering
\usepackage{iftex}
\ifPDFTeX
  \usepackage[T1]{fontenc}
  \usepackage[utf8]{inputenc}
  \usepackage{textcomp} % provide euro and other symbols
\else % if luatex or xetex
  \usepackage{unicode-math} % this also loads fontspec
  \defaultfontfeatures{Scale=MatchLowercase}
  \defaultfontfeatures[\rmfamily]{Ligatures=TeX,Scale=1}
\fi
\usepackage{lmodern}
\ifPDFTeX\else
  % xetex/luatex font selection
\fi
% Use upquote if available, for straight quotes in verbatim environments
\IfFileExists{upquote.sty}{\usepackage{upquote}}{}
\IfFileExists{microtype.sty}{% use microtype if available
  \usepackage[]{microtype}
  \UseMicrotypeSet[protrusion]{basicmath} % disable protrusion for tt fonts
}{}
\makeatletter
\@ifundefined{KOMAClassName}{% if non-KOMA class
  \IfFileExists{parskip.sty}{%
    \usepackage{parskip}
  }{% else
    \setlength{\parindent}{0pt}
    \setlength{\parskip}{6pt plus 2pt minus 1pt}}
}{% if KOMA class
  \KOMAoptions{parskip=half}}
\makeatother
\usepackage{color}
\usepackage{fancyvrb}
\newcommand{\VerbBar}{|}
\newcommand{\VERB}{\Verb[commandchars=\\\{\}]}
\DefineVerbatimEnvironment{Highlighting}{Verbatim}{commandchars=\\\{\}}
% Add ',fontsize=\small' for more characters per line
\usepackage{framed}
\definecolor{shadecolor}{RGB}{248,248,248}
\newenvironment{Shaded}{\begin{snugshade}}{\end{snugshade}}
\newcommand{\AlertTok}[1]{\textcolor[rgb]{0.94,0.16,0.16}{#1}}
\newcommand{\AnnotationTok}[1]{\textcolor[rgb]{0.56,0.35,0.01}{\textbf{\textit{#1}}}}
\newcommand{\AttributeTok}[1]{\textcolor[rgb]{0.13,0.29,0.53}{#1}}
\newcommand{\BaseNTok}[1]{\textcolor[rgb]{0.00,0.00,0.81}{#1}}
\newcommand{\BuiltInTok}[1]{#1}
\newcommand{\CharTok}[1]{\textcolor[rgb]{0.31,0.60,0.02}{#1}}
\newcommand{\CommentTok}[1]{\textcolor[rgb]{0.56,0.35,0.01}{\textit{#1}}}
\newcommand{\CommentVarTok}[1]{\textcolor[rgb]{0.56,0.35,0.01}{\textbf{\textit{#1}}}}
\newcommand{\ConstantTok}[1]{\textcolor[rgb]{0.56,0.35,0.01}{#1}}
\newcommand{\ControlFlowTok}[1]{\textcolor[rgb]{0.13,0.29,0.53}{\textbf{#1}}}
\newcommand{\DataTypeTok}[1]{\textcolor[rgb]{0.13,0.29,0.53}{#1}}
\newcommand{\DecValTok}[1]{\textcolor[rgb]{0.00,0.00,0.81}{#1}}
\newcommand{\DocumentationTok}[1]{\textcolor[rgb]{0.56,0.35,0.01}{\textbf{\textit{#1}}}}
\newcommand{\ErrorTok}[1]{\textcolor[rgb]{0.64,0.00,0.00}{\textbf{#1}}}
\newcommand{\ExtensionTok}[1]{#1}
\newcommand{\FloatTok}[1]{\textcolor[rgb]{0.00,0.00,0.81}{#1}}
\newcommand{\FunctionTok}[1]{\textcolor[rgb]{0.13,0.29,0.53}{\textbf{#1}}}
\newcommand{\ImportTok}[1]{#1}
\newcommand{\InformationTok}[1]{\textcolor[rgb]{0.56,0.35,0.01}{\textbf{\textit{#1}}}}
\newcommand{\KeywordTok}[1]{\textcolor[rgb]{0.13,0.29,0.53}{\textbf{#1}}}
\newcommand{\NormalTok}[1]{#1}
\newcommand{\OperatorTok}[1]{\textcolor[rgb]{0.81,0.36,0.00}{\textbf{#1}}}
\newcommand{\OtherTok}[1]{\textcolor[rgb]{0.56,0.35,0.01}{#1}}
\newcommand{\PreprocessorTok}[1]{\textcolor[rgb]{0.56,0.35,0.01}{\textit{#1}}}
\newcommand{\RegionMarkerTok}[1]{#1}
\newcommand{\SpecialCharTok}[1]{\textcolor[rgb]{0.81,0.36,0.00}{\textbf{#1}}}
\newcommand{\SpecialStringTok}[1]{\textcolor[rgb]{0.31,0.60,0.02}{#1}}
\newcommand{\StringTok}[1]{\textcolor[rgb]{0.31,0.60,0.02}{#1}}
\newcommand{\VariableTok}[1]{\textcolor[rgb]{0.00,0.00,0.00}{#1}}
\newcommand{\VerbatimStringTok}[1]{\textcolor[rgb]{0.31,0.60,0.02}{#1}}
\newcommand{\WarningTok}[1]{\textcolor[rgb]{0.56,0.35,0.01}{\textbf{\textit{#1}}}}
\usepackage{graphicx}
\makeatletter
\newsavebox\pandoc@box
\newcommand*\pandocbounded[1]{% scales image to fit in text height/width
  \sbox\pandoc@box{#1}%
  \Gscale@div\@tempa{\textheight}{\dimexpr\ht\pandoc@box+\dp\pandoc@box\relax}%
  \Gscale@div\@tempb{\linewidth}{\wd\pandoc@box}%
  \ifdim\@tempb\p@<\@tempa\p@\let\@tempa\@tempb\fi% select the smaller of both
  \ifdim\@tempa\p@<\p@\scalebox{\@tempa}{\usebox\pandoc@box}%
  \else\usebox{\pandoc@box}%
  \fi%
}
% Set default figure placement to htbp
\def\fps@figure{htbp}
\makeatother
\setlength{\emergencystretch}{3em} % prevent overfull lines
\providecommand{\tightlist}{%
  \setlength{\itemsep}{0pt}\setlength{\parskip}{0pt}}
\usepackage{bookmark}
\IfFileExists{xurl.sty}{\usepackage{xurl}}{} % add URL line breaks if available
\urlstyle{same}
\hypersetup{
  pdftitle={Music Motor Manuscript},
  hidelinks,
  pdfcreator={LaTeX via pandoc}}

\title{Music Motor Manuscript}
\author{}
\date{\vspace{-2.5em}2026-03-12}

\begin{document}
\maketitle

\subsection{Abstract}\label{abstract}

Prior research suggests that music-evoked emotions may drive movement,
thereby reinforcing emotional experience through engagement of the motor
system. This relationship has been linked to activity in the primary
motor cortex (M1); however, its causal role remains unclear. The present
within-subjects study examined whether inhibiting M1 using continuous
theta-burst stimulation alters perceived valence and arousal ratings of
music. Contrary to our hypothesis, no significant effects of M1
inhibition on valence or arousal ratings were observed relative to sham
stimulation. Exploratory analyses examining musical sophistication, as
measured by the Goldsmith Musical Sophistication Index (Gold-MSI),
revealed that arousal ratings were significantly associated with the
Active Engagement and Emotional Engagement subscales. In contrast, no
significant relationships were observed between Gold-MSI scores and
valence ratings. Gold-MSI scores did not moderate the effects of
stimulation condition. These findings suggest that M1 may not play a
direct causal role in emotional responses to music, highlighting the
need to investigate other motor-related regions.

Keywords: primary motor cortex (M1); continuous theta-burst stimulation
(cTBS); valence; arousal; music.

\begin{verbatim}
\end{verbatim}

\subsection{Introduction}\label{introduction}

Emotion and movement are tightly interconnected, as physical action is
not only an expression of emotion but also a contributor to its
production (Shafir et al., 2013). This idea is rooted in the
James--Lange theory, which proposes that emotions arise as a response to
bodily changes (e.g., fear emerges from physiological reactions such as
trembling and increased heart rate). Damasio's somatic marker hypothesis
further extends this framework by suggesting that the brain stores
bodily states as memories and links them to emotional experiences. As a
result, when individuals perform, observe, or imagine a movement, the
brain reactivates these bodily states to generate emotion in the current
context. Consistent with these theories, Shafir et al.~(2013)
demonstrated that executing, observing, and imagining body movements
associated with specific emotions (e.g., happiness, sadness, fear)
increased the corresponding affective state. The cross-modal bias theory
further extends these frameworks by proposing that emotional information
from one sensory modality, such as music, influences how emotions are
experienced in another modality, such as body movement (Christensen et
al., 2014). Christensen et al.~(2014) demonstrated that emotional
experience can be characterized along two dimensions: valence, which
refers to how positive or negative an emotion is, and arousal, which
refers to the level of energy. Their results showed that subjective
emotional ratings were primarily driven by valence; ergo, music
influenced whether dance movements were perceived as more positive or
negative. In contrast, physiological responses were modulated by arousal
congruency, with the strongest bodily responses occurring when
low-arousal movements were paired with sad music and high-arousal
movements were paired with happy music. Yet, despite growing evidence of
an association between movement and emotion, only recently has research
begun to employ noninvasive neurophysiological techniques to investigate
this correlation. Doruk et al.~(2016) helped address this gap in the
literature by providing evidence that specific neurophysiological
markers of the motor system are associated with post-stroke emotional
disturbances, such as anxiety and depression. Using
electroencephalography (EEG) and transcranial magnetic stimulation
(TMS), the study demonstrated that emotional disturbance was not solely
attributable to limbic dysfunction, but also reflected altered motor
system dynamics. EEG revealed reduced interhemispheric coherence and
frontal--central connectivity, while TMS indicated increased
intracortical facilitation and higher motor thresholds in the unlesioned
hemisphere. These neurophysiological markers were not associated with
motor impairment severity; however, they were associated with mood and
emotional control. Collectively, these findings indicate that motor and
emotional disturbances following stroke are supported by distinct yet
interacting neural mechanisms, rather than emotional disturbance being a
direct consequence of motor impairment. Selosse et al.~(2025) sought to
further evaluate the relationship between movement and emotion by
isolating components of the motor system, specifically examining the
relationship between the primary motor cortex (M1) and emotion. The
researchers used functional magnetic resonance imaging (fMRI) to examine
emotional prosody production while participants produced pseudowords in
angry, happy, and neutral tones under three conditions: normal speech,
whispering, and silent articulation. The authors found greater M1
activation during normally voiced emotional speech compared to
whispering or silent articulation. This pattern suggests that emotional
prosody relies on a broader neural network that includes motor regions,
with M1 supporting the execution of the vocal movements required to
express emotion. Thus, emotional voice production can be understood as
an embodied process in which M1 contributes to the physical expression
of emotion through vocal tract movement and speech. Xia et al.~(2021)
further examined the relationship between the M1 and emotion by
distinguishing the associations between emotional processing and
activity in the left and right M1. The study employed TMS to record
motor-evoked potentials (MEPs) from both the left and right M1 while
right-handed participants were instructed to move a manikin toward
positively valenced images and away from negatively valenced images
using keyboard presses. The results showed that the left M1 exhibited
larger MEP amplitudes, indicating increased excitability, when
participants performed compatible behaviors (i.e., approaching positive
stimuli and avoiding negative stimuli). In contrast, the right M1 showed
reduced excitability during these compatible behaviors. These effects
emerged prior to physical movement; therefore, during stimulus
processing. Together, these findings suggest that emotional valence
modulates motor preparation, supporting a functional association between
emotion processing and contralateral M1 activity. Despite growing
evidence linking M1 activity to emotional processing, most research has
established only correlational relationships. Establishing causality
requires noninvasive neuromodulation techniques, such as continuous
theta-burst stimulation (cTBS), which can inhibit cortical regions to
assess their functional contributions. Prior work has demonstrated that
cTBS applied to M1 can be used to test for causal involvement by
examining behavioral consequences following inhibition (e.g., Kraeutner
et al., 2017). However, little research has directly investigated
whether inhibiting M1 alters self-reported emotions, leaving the causal
role of M1 in emotion largely unresolved. Hence, the present study aims
to address this gap by using cTBS to inhibit the contralateral M1 and
examining its effects on emotion production in response to pre-rated
music clips varying in valence and arousal. In a within-subjects design,
right-handed participants completed two sessions: active cTBS
(experimental) and sham (control). After stimulation, participants
listened to instrumental music clips and rated each clip's valence and
arousal. Participant musical sophistication was also measured using a
questionnaire. We hypothesize that M1 inhibition will yield more neutral
valence and arousal ratings than sham, reflecting reduced M1 involvement
in emotion processing. Additionally, if participants possess greater
musical sophistication, they may show more polarized ratings under M1
inhibition than those with lower musical sophistication, suggesting
alternative emotion-processing mechanisms less dependent on M1.

\begin{verbatim}
\end{verbatim}

\subsection{Methods}\label{methods}

\textbf{Participants}

Study participants (N = 29) were recruited through the SONA online
platform and poster advertisement. Eligibility was restricted to
right-handed individuals who were able to understand English and had
normal or corrected-to-normal vision. Participants were excluded if they
self-reported any neurological conditions or impairments that would
affect motor function, or if they had any contraindications to TMS as
identified by a standard screening form (Rossi et al., 2009). For the
final analysis, N = 22 university students were included, and 14
identified as female. Participants had a mean age of 22.73 years (SD =
4.7, range = 19--38) and a mean General Musical Sophistication score of
74.68 (SD = 16.45, range = 35--114). Each participant received either
\$10 CAD per session or SONA course credit (2 credits for Session 1 and
1.5 credits for Session 2). Data collection took place between July 15,
2025, and February 26, 2026.

\textbf{Materials}

\textbf{Transcranial Magnetic Stimulation (TMS) Screening Form}

To ensure participant eligibility and safety during transcranial
magnetic stimulation (TMS), the TMS Screening Form was administered
(Rossi et al., 2009). The questionnaire includes 10 binary items
(yes/no) assessing medical history; a ``yes'' response to any item
results in exclusion. Participants were also required to report current
medications, the presence of metal in the body, and any prior experience
with TMS or magnetic resonance imaging (MRI). Any medications, implanted
metal, or history indicating a contraindication or adverse response to
TMS or MRI led to exclusion.

\textbf{Kinaesthetic and Visual Imagery Questionnaire (KVIQ)}

The Kinaesthetic and Visual Imagery Questionnaire (KVIQ) was used to
assess participants' ability to perform kinaesthetic and visual motor
imagery (Malouin et al., 2007). Participants physically performed five
body movements and then imagined each movement using visual imagery,
followed by kinaesthetic imagery. Participants rated the clarity and
intensity of the imagined movement on a 5-point scale, with 1 indicating
the weakest imagery and 5 indicating the strongest.

\textbf{Goldsmith Musical Sophistication Index (Gold-MSI)}

To assess musical sophistication, participants completed the Goldsmiths
Musical Sophistication Index (Gold-MSI; Müllensiefen et al., 2014). This
questionnaire evaluates six dimensions: active engagement, perceptual
ability, musical training, emotional engagement, singing ability, and
general musical sophistication. Participants responded to items using a
7-point Likert scale, with 1 indicating strong disagreement and 7
indicating strong agreement. The questionnaire was also used to collect
descriptive information, including participants' age, sex, and education
level.

\textbf{Transcranial Magnetic Stimulation (TMS)}

Participants underwent transcranial magnetic stimulation (TMS)-- a
non-invasive technique that uses brief magnetic pulses to generate
electrical currents in the brain and will be used in accordance with
well-established procedures (Rossi et al.~2009; Huang et al.~2005;
Wischnewski et al.~2015). Stimulation output intensity has been
individualized using established protocols (Kleim et al., 2007; Rossi et
al., 2009). Neuronavigation (Brainsight 2™; Rogue Research Inc.)
provided focal, real-time localization of the TMS coil with millimetric
accuracy by optically tracking coil position relative to digitized
anatomical landmarks co-registered to a template brain. This system
enabled precise targeting of cortical regions. To determine stimulation
intensity and the motor ``hotspot,'' resting motor threshold (RMT) was
assessed as an index of cortical excitability. RMT is defined as the
minimum stimulator output that elicits visible muscle twitches and
motor-evoked potentials ≥50 μV. This was recorded using surface
electromyography (EMG) over the left first dorsal interosseous (FDI) in
at least 5 of 10 trials. Baseline motor excitability was assessed by
recording 16 motor evoked potentials (MEPs) at rest, 16 during
kinaesthetic motor imagery (KMI), and 16 during visual motor imagery
(VMI) of a simple finger abduction. TMS was delivered at 120\% of
resting motor threshold (RMT) (Bruton et al., 2020; Stinear et al.,
2007; Wright et al., 2014). Participants received continuous theta-burst
stimulation (cTBS) over the contralateral M1 in accordance with
established guidelines (Huang et al., 2005; Rossi et al., 2009, 2020;
Wischnewski et al., 2015). The cTBS protocol consisted of bursts of
three pulses at 50 Hz, repeated every 200 ms for 40 s (600 pulses) at
80\% RMT, delivered using a Magstim SuperRapid2 Plus-1 with a cooled
figure-of-eight coil. Sham cTBS was delivered to the vertex at 20\% RMT
for an equivalent duration using a validated protocol that mimics the
sensory experience of TMS without producing physiological effects
(Kraeutner et al., 2016, 2017). All stimulation parameters were within
established safety standards. Following cTBS or sham stimulation, MEPs
were recorded under the same conditions to confirm M1 inhibition.

\textbf{Active Listening Task}

Participants' emotion production was assessed by completing an active
listening task, which followed previously reported procedures (Eerola \&
Vuoskoski, 2011). Participants were randomly assigned to listen to one
of two sets of 32 pre-rated instrumental music clips, with clip order
randomized across sessions and participants. The clips were balanced for
valence and arousal and were selected from an established set of 110
clips (openly available at \url{https://osf.io/p6vkg/wiki/home/}).
During the task, participants were seated in a chair with their dominant
hand resting on their lap and their non-dominant hand placed on a
keypad. They listened to the clips through headphones and, after each
clip, used the keypad to rate their emotional responses on 9-point
valence and arousal scales. Procedure Participants were informed that
participation was voluntary and that they could withdraw from the study
at any time without penalty or loss of compensation. All participants
completed a TMS safety screening form to confirm eligibility and
provided written informed consent before participation. See Figure 1 for
a schematic overview of the study procedure. The study consisted of two
sessions separated by a minimum of 24 hours. In Session 1, participants
completed the KVIQ, the Motor Experience Questionnaire, and the
Gold-MSI. Participants were then seated comfortably, and surface EMG
electrodes were placed on the left FDI to record muscle activity. Using
neuronavigation, the contralateral M1 was localized, and RMT was
determined. Participants subsequently received 16 single-pulse TMS
stimulations at rest, 16 during KMI, and 16 during VMI. The motor
imagery task consisted of imagining abduction of the left index finger.
Imagery trials were cued using the software PsychoPy; participants were
instructed to close their eyes and begin imagining upon hearing the
auditory cue ``Go'' and to stop imagining upon hearing ``Rest.'' Each
imagery period lasted 10 seconds, during which two TMS pulses were
delivered. Participants were permitted to rest between trials and
initiated each subsequent trial by pressing a keyboard key. Following
baseline stimulation, participants received either active or sham cTBS.
Earplugs were provided to minimize auditory discomfort. All participants
received both active and sham cTBS across two sessions, with the order
randomized. Participants were informed that they would receive both
conditions, but were blinded to the condition administered in each
session. After cTBS, participants were instructed to refrain from moving
their left hand for 40 minutes. A 10-minute rest period followed cTBS,
after which participants again received 16 TMS pulses at rest, 16 during
KMI, and 16 during VMI. The order of KMI and VMI conditions was
counterbalanced across sessions. Participants then completed an active
listening task. Using a PsychoPy program, participants listened to 32
instrumental music clips through headphones. After each clip, they rated
perceived valence and arousal on a 9-point scale using keyboard
responses. The order of valence and arousal ratings was counterbalanced
across sessions. All post-cTBS tasks were completed within 40 minutes,
and no data were collected beyond this period. At the conclusion of each
session, participants received compensation in the form of cash or SONA
research credit. See Figure 1 for a visual of the study methodology.

\textbf{Statistical analyses}

All statistical analyses were conducted using RStudio. Participants were
randomly assigned to receive cTBS in session 1 or session 2. The order
of motor imagery task performance---KMI and VMI---as well as the rating
of valence and arousal, was counterbalanced across sessions. Each
participant contributed eight observations per session, corresponding to
ratings of the four pre-defined categories of valence and arousal: high
valence--low arousal, low valence--high arousal, high valence--high
arousal, and low valence--low arousal. Participants provided ratings in
both the sham control and active cTBS conditions. Univariate analysis of
variance (ANOVA) were conducted to assess the effects of active versus
sham cTBS on valence and arousal ratings across the four clip
categories. Linear mixed-effects models (LMEs) were conducted to examine
whether Gold-MSI ratings mediated the relationship between M1 inhibition
and perceived valence and arousal ratings. Music clips were selected
from a validated set of 110 clips. Four categories were created: 16
clips with the largest difference between high valence and low arousal,
16 clips with the largest difference between low valence and high
arousal, 16 clips with the highest valence and arousal, and 16 clips
with the lowest valence and arousal. These four groups were then
randomly divided into two sets of eight clips per category, yielding two
sets of 32 music clips. Participants were randomly assigned to listen to
one of the two sets. The presentation order of the clips was randomized
within each participant's session and across participants.

\begin{verbatim}
\end{verbatim}

\subsection{Results}\label{results}

Following the criteria above, seven participants were excluded due to
inability to establish a reliable motor threshold (n = 4), issues during
the session (n = 1), use of medication contraindicated for TMS (n = 1),
and participant failure to attend the session (n = 1). Hence, data from
22 participants were included in the final analysis.

\textbf{Active Listening Task}

To test the hypothesis that inhibition of M1 would neutralize perceived
valence and arousal ratings, two separate univariate analyses of
variance (ANOVAs) were conducted. For both mean valence and mean arousal
ratings, there was no significant main effect of stimulation condition
(see Figures 2 \& 3). However, a significant main effect of stimulus
type was observed for both valence (p \textless{} .001) and arousal (p
\textless{} .001) using Greenhouse--Geisser corrections (see Figures 2
\& 3). This indicates that ratings differed across the four stimulus
categories. Effect size analyses further indicated that the stimulation
condition had minimal influence on both valence and arousal ratings. For
valence, Cohen's d values were small or negligible across stimulus
categories: high valence--low arousal (d = 0.37), indicating a
small-to-moderate effect with higher ratings in the active condition;
low valence--low arousal (d = 0.23), indicating a small effect with
higher ratings in the active condition; low valence--high arousal (d =
−0.15), indicating a very small effect with lower ratings in the active
condition; and high valence--high arousal (d = 0.01), indicating no
meaningful effect. Similarly, for arousal ratings, effect sizes were
small or negligible: high valence--low arousal (d = −0.23), indicating a
small effect with lower ratings in the active condition; low
valence--high arousal (d = 0.15), indicating a very small effect with
higher ratings in the active condition; low valence--low arousal (d =
0.07), indicating no meaningful effect; and high valence--high arousal
(d = 0.02), also indicating no meaningful effect.

\textbf{General Musical Sophistication}

To examine whether musical sophistication was associated with perceived
valence and arousal ratings across stimulus types, linear mixed-effects
models (LMEs) were conducted, with participant included as a random
intercept. Gold-MSI scores (Active Engagement, Perceptual Ability,
Musical Training, Emotional Engagement, Singing Ability, and General
Musical Sophistication) were entered as fixed effects, along with
stimulus type. The model predicting mean valence ratings revealed no
significant main effects of Gold-MSI scores on any subscale, indicating
that individual differences in musical sophistication did not
significantly predict mean valence ratings across stimulus categories.
The model predicting mean arousal ratings revealed significant effects
of the Gold-MSI subscales Active Engagement (p = .004) and Emotional
Engagement (p = .005; see Figures 4 \& 5). In contrast, the remaining
subscales (Perceptual Ability, Musical Training, Singing Ability, and
General Musical Sophistication) were not significantly associated with
mean arousal ratings. Additional LMEs examining both mean valence and
arousal ratings across stimulus type and stimulation condition revealed
no significant main effects involving Gold-MSI scores across the six
subscales.

\begin{verbatim}
\end{verbatim}

\subsection{Discussion}\label{discussion}

We hypothesized that inhibition of M1 would lead to more neutral valence
and arousal ratings compared to sham stimulation, reflecting a reduced
role of M1 in emotion processing. Additionally, we predicted that
individuals with greater musical sophistication would exhibit more
polarized emotional ratings under M1 inhibition, suggesting reliance on
alternative emotion-processing mechanisms. Contrary to these
predictions, stimulation condition did not significantly influence
either valence or arousal ratings, and effect size analyses indicated
only minimal differences between active and sham conditions. Instead,
emotional ratings were primarily driven by stimulus type, with clear
differences observed across the four stimulus categories. Furthermore,
musical sophistication did not moderate the effects of M1 inhibition on
emotional ratings. Although exploratory analyses revealed that Active
Engagement and Emotional Engagement were associated with higher arousal
ratings, no such relationships were observed for valence. Together,
these findings do not support a causal role of M1 in shaping emotional
responses to music, nor do they indicate that musical sophistication
alters this relationship. One possible explanation for the absence of M1
effects is that the active listening task may not have strongly engaged
motor-related processes (Lee \& Siegle, 2012). Prior research suggests
that different types of emotional evaluations recruit distinct neural
networks. For example, stimulus-focused evaluations, such as rating the
emotional content of a stimulus, primarily engage sensory processing
regions and prefrontal areas involved in evaluation, including the
ventrolateral prefrontal cortex. In contrast, evaluating one's own
emotional state recruits regions such as the insula and anterior
cingulate cortex, while evaluating others' emotions engages
social-cognitive regions such as the superior temporal sulcus and
temporoparietal junction. Since participants in the present study rated
the perceived valence and arousal of musical stimuli, their responses
may have relied more on perceptual and evaluative networks than on motor
system involvement. This may help explain why inhibiting M1 did not
significantly alter emotional ratings. The lack of M1 involvement in the
present study may also be explained by the task's limited social and
interactive context (Antonelli et al., 2025). Inter-neural synchrony
(INS), which reflects coordinated brain activity across individuals, may
emerge most strongly during real-time, reciprocal social interactions
where individuals continuously exchange emotional and behavioral cues.
In contrast, when participants engage with non-social or
computer-mediated stimuli, such as rating music on a device, the
interaction lacks the dynamic feedback loop necessary to support strong
interpersonal coupling. This, in turn, may reduce INS. Since motor
regions such as M1 are often implicated in action coordination and
interpersonal synchrony, reduced INS in a non-interactive setting may
limit recruitment of these regions. Consequently, the absence of a
meaningful social context in this study may have attenuated M1
involvement; thereby, inhibitory stimulation did not significantly alter
participants' emotional ratings of the music. Another possible
explanation for why inhibiting the M1 did not alter participants'
valence or arousal ratings of music clips is that M1 may not be central
in generating affective evaluations of music (Putkinen et al., 2021).
Although the motor cortex contains emotion-specific activity patterns
and supports embodied aspects of music perception (e.g., movement),
fundamental affective dimensions such as valence and arousal are
primarily mediated by auditory cortical processing and by broader limbic
and paralimbic systems. As a result, any potential influence of musical
sophistication on M1-related emotional processing may not be detectable
because there are no observable differences in valence or arousal
experiences relative to sham. Although exploratory analyses showed that
Active Engagement and Emotional Engagement were associated with higher
arousal ratings, no such relationships were observed for valence. One
possible explanation is that participants who are more emotionally and
actively engaged with music may experience stronger empathic or embodied
responses during the active listening task (Chapin et al., 2010). Music
may evoke emotion by prompting participants to internally simulate the
tempo and its expression, making the experience feel more intense.
Individuals who frequently engage with music or are more emotionally
attuned may be more sensitive to these expressive cues, leading to
heightened arousal. In contrast, valence appears to be more subjective
and less consistently tied to these embodied processes, which may
explain why Active and Emotional Engagement did not influence whether
emotions were perceived as positive or negative. Future research should
move beyond numerical valence and arousal scales, as they are limited in
their ability to represent emotion and can be interpreted differently by
participants, introducing variability in the data. Future studies could
include discrete emotion categories such as happiness, sadness, fear,
anger, and tenderness, alongside dimensional measures, to better capture
the specific emotional aspects of music (Eerola \& Vuoskoski, 2011).
Additionally, distinguishing among types of arousal, such as energy and
tension arousal, may provide more detail about how emotions are
experienced. Including measures such as preference and perceived beauty
would also help separate emotional valence from aesthetic liking. Since
prior research suggests that dimensional and discrete models capture
different aspects of emotion and may rely on different processes,
combining these approaches could provide a more accurate and nuanced
understanding of emotional responses to music.

\textbf{Limitations}

Several limitations should be noted. First, the sample size may have
been too small to detect subtle effects. We also did not recruit
participants across the full range of musical sophistication, which
likely reduced variability and made it harder to observe any moderating
effects. Using only Western music clips may limit emotional responses
and reduce the generalizability of the findings across cultures. In
addition, the laboratory setting was highly controlled and may not
reflect real-world musical experiences. Our sample consisted mainly of
university students in their early twenties, which limits
generalizability to other age groups. We did not analyze Brainsight data
to confirm accurate targeting of M1. We also did not verify that
cortical inhibition was maintained after the music listening task. These
factors may have reduced our ability to detect meaningful effects.

\textbf{Conclusion}

In conclusion, our findings challenge prior research suggesting that M1
plays a direct role in emotion production, as we observed no causal
relationship. Emotional ratings were not influenced by musical
sophistication under M1 inhibition. However, participants with higher
Active and Emotional Engagement with music showed greater sensitivity to
arousal. These results help clarify which neural systems are involved in
emotional processing. This is important for guiding future research, as
it suggests that targeting M1 alone may not be effective for influencing
emotion. Instead, it highlights the need to focus on broader networks
when developing interventions for motor rehabilitation and the treatment
of affective disorders.

\subsection{References}\label{references}

Shafir, T., Taylor, S. F., Atkinson, A. P., Langenecker, S. A., \&
Zubieta, J. K. (2013). Emotion~

regulation through execution, observation, and imagery of emotional
movements. Brain~

and Cognition, 82(2), 219-227.
\url{https://doi.org/10.1016/j.bandc.2013.03.001}

Christensen, J. F., Gaigg, S. B., Gomila, A., Oke, P., \& Calvo-Merino,
B. (2014). Enhancing~

emotional experiences to dance through music: the role of valence and
arousal in the~

cross-modal bias. Frontiers In Human Neuroscience, 8, 757.~

\url{https://doi.org/10.3389/fnhum.2014.00757}

\subsection{Plots}\label{plots}

\begin{Shaded}
\begin{Highlighting}[]
\NormalTok{data }\OtherTok{\textless{}{-}} \FunctionTok{read.csv}\NormalTok{(here}\SpecialCharTok{::}\FunctionTok{here}\NormalTok{(}\StringTok{"data"}\NormalTok{, }\StringTok{"processed"}\NormalTok{, }\StringTok{"music\_task\_data\_participant{-}level.csv"}\NormalTok{))}

\FunctionTok{library}\NormalTok{(here)}
\end{Highlighting}
\end{Shaded}

\begin{verbatim}
## Warning: package 'here' was built under R version 4.4.3
\end{verbatim}

\begin{verbatim}
## here() starts at C:/Users/knopf/OneDrive/Documents/school_2025-2026/Biol430G/BIOL430G_mini-project
\end{verbatim}

\begin{Shaded}
\begin{Highlighting}[]
\FunctionTok{library}\NormalTok{(ez)}
\end{Highlighting}
\end{Shaded}

\begin{verbatim}
## Warning: package 'ez' was built under R version 4.4.3
\end{verbatim}

\begin{Shaded}
\begin{Highlighting}[]
\FunctionTok{library}\NormalTok{(dplyr)}
\end{Highlighting}
\end{Shaded}

\begin{verbatim}
## 
## Attaching package: 'dplyr'
\end{verbatim}

\begin{verbatim}
## The following objects are masked from 'package:stats':
## 
##     filter, lag
\end{verbatim}

\begin{verbatim}
## The following objects are masked from 'package:base':
## 
##     intersect, setdiff, setequal, union
\end{verbatim}

\begin{Shaded}
\begin{Highlighting}[]
\FunctionTok{library}\NormalTok{(tidyr)}
\FunctionTok{library}\NormalTok{(ggplot2)}
\end{Highlighting}
\end{Shaded}

\begin{verbatim}
## Warning: package 'ggplot2' was built under R version 4.4.3
\end{verbatim}

\begin{Shaded}
\begin{Highlighting}[]
\FunctionTok{library}\NormalTok{(lme4)}
\end{Highlighting}
\end{Shaded}

\begin{verbatim}
## Loading required package: Matrix
\end{verbatim}

\begin{verbatim}
## 
## Attaching package: 'Matrix'
\end{verbatim}

\begin{verbatim}
## The following objects are masked from 'package:tidyr':
## 
##     expand, pack, unpack
\end{verbatim}

\begin{Shaded}
\begin{Highlighting}[]
\FunctionTok{library}\NormalTok{(lmerTest)}
\end{Highlighting}
\end{Shaded}

\begin{verbatim}
## Warning: package 'lmerTest' was built under R version 4.4.3
\end{verbatim}

\begin{verbatim}
## 
## Attaching package: 'lmerTest'
\end{verbatim}

\begin{verbatim}
## The following object is masked from 'package:lme4':
## 
##     lmer
\end{verbatim}

\begin{verbatim}
## The following object is masked from 'package:stats':
## 
##     step
\end{verbatim}

\begin{Shaded}
\begin{Highlighting}[]
\FunctionTok{library}\NormalTok{(effectsize)}
\end{Highlighting}
\end{Shaded}

\begin{verbatim}
## Warning: package 'effectsize' was built under R version 4.4.3
\end{verbatim}

\begin{Shaded}
\begin{Highlighting}[]
\FunctionTok{library}\NormalTok{(renv)}
\end{Highlighting}
\end{Shaded}

\begin{verbatim}
## Warning: package 'renv' was built under R version 4.4.3
\end{verbatim}

\begin{verbatim}
## 
## Attaching package: 'renv'
\end{verbatim}

\begin{verbatim}
## The following object is masked from 'package:Matrix':
## 
##     update
\end{verbatim}

\begin{verbatim}
## The following objects are masked from 'package:stats':
## 
##     embed, update
\end{verbatim}

\begin{verbatim}
## The following objects are masked from 'package:utils':
## 
##     history, upgrade
\end{verbatim}

\begin{verbatim}
## The following objects are masked from 'package:base':
## 
##     autoload, load, remove, use
\end{verbatim}

\begin{Shaded}
\begin{Highlighting}[]
\NormalTok{renv}\SpecialCharTok{::}\FunctionTok{snapshot}\NormalTok{()}
\end{Highlighting}
\end{Shaded}

\begin{verbatim}
## The following package(s) will be updated in the lockfile:
## 
## # CRAN -----------------------------------------------------------------------
## - abind            [* -> 1.4-8]
## - backports        [* -> 1.5.0]
## - base64enc        [* -> 0.1-3]
## - bayestestR       [* -> 0.17.0]
## - boot             [* -> 1.3-31]
## - broom            [* -> 1.0.7]
## - bslib            [* -> 0.9.0]
## - cachem           [* -> 1.1.0]
## - car              [* -> 3.1-3]
## - carData          [* -> 3.0-5]
## - cli              [* -> 3.6.3]
## - colorspace       [* -> 2.1-1]
## - cowplot          [* -> 1.1.3]
## - cpp11            [* -> 0.5.2]
## - datawizard       [* -> 1.3.0]
## - Deriv            [* -> 4.1.6]
## - digest           [* -> 0.6.37]
## - doBy             [* -> 4.6.25]
## - dplyr            [* -> 1.1.4]
## - effectsize       [* -> 1.0.2]
## - evaluate         [* -> 1.0.3]
## - ez               [* -> 4.5-0]
## - fansi            [* -> 1.0.6]
## - farver           [* -> 2.1.2]
## - fastmap          [* -> 1.2.0]
## - fontawesome      [* -> 0.5.3]
## - Formula          [* -> 1.2-5]
## - fs               [* -> 1.6.5]
## - generics         [* -> 0.1.3]
## - ggplot2          [* -> 3.5.2]
## - glue             [* -> 1.8.0]
## - gtable           [* -> 0.3.6]
## - here             [* -> 1.0.2]
## - highr            [* -> 0.11]
## - htmltools        [* -> 0.5.8.1]
## - insight          [* -> 1.4.6]
## - isoband          [* -> 0.2.7]
## - jquerylib        [* -> 0.1.4]
## - jsonlite         [* -> 1.8.9]
## - knitr            [* -> 1.49]
## - labeling         [* -> 0.4.3]
## - lattice          [* -> 0.22-6]
## - lifecycle        [* -> 1.0.4]
## - lme4             [* -> 1.1-36]
## - lmerTest         [* -> 3.1-3]
## - magrittr         [* -> 2.0.3]
## - MASS             [* -> 7.3-65]
## - Matrix           [* -> 1.7-1]
## - MatrixModels     [* -> 0.5-3]
## - memoise          [* -> 2.0.1]
## - mgcv             [* -> 1.9-1]
## - microbenchmark   [* -> 1.5.0]
## - mime             [* -> 0.12]
## - minqa            [* -> 1.2.8]
## - modelr           [* -> 0.1.11]
## - munsell          [* -> 0.5.1]
## - nlme             [* -> 3.1-167]
## - nloptr           [* -> 2.1.1]
## - nnet             [* -> 7.3-19]
## - numDeriv         [* -> 2016.8-1.1]
## - parameters       [* -> 0.28.3]
## - pbkrtest         [* -> 0.5.3]
## - performance      [* -> 0.16.0]
## - pillar           [* -> 1.10.1]
## - pkgconfig        [* -> 2.0.3]
## - plyr             [* -> 1.8.9]
## - purrr            [* -> 1.0.4]
## - quantreg         [* -> 6.00]
## - R6               [* -> 2.6.1]
## - rappdirs         [* -> 0.3.3]
## - rbibutils        [* -> 2.3]
## - RColorBrewer     [* -> 1.1-3]
## - Rcpp             [* -> 1.0.14]
## - RcppEigen        [* -> 0.3.4.0.2]
## - Rdpack           [* -> 2.6.2]
## - reformulas       [* -> 0.4.0]
## - renv             [* -> 1.1.8]
## - reshape2         [* -> 1.4.4]
## - rlang            [* -> 1.1.5]
## - rmarkdown        [* -> 2.30]
## - rprojroot        [* -> 2.1.1]
## - sass             [* -> 0.4.9]
## - scales           [* -> 1.3.0]
## - SparseM          [* -> 1.84-2]
## - stringi          [* -> 1.8.4]
## - stringr          [* -> 1.5.1]
## - survival         [* -> 3.7-0]
## - tibble           [* -> 3.2.1]
## - tidyr            [* -> 1.3.1]
## - tidyselect       [* -> 1.2.1]
## - tinytex          [* -> 0.54]
## - utf8             [* -> 1.2.4]
## - vctrs            [* -> 0.6.5]
## - viridisLite      [* -> 0.4.2]
## - withr            [* -> 3.0.2]
## - xfun             [* -> 0.50]
## - yaml             [* -> 2.3.10]
## 
## The version of R recorded in the lockfile will be updated:
## - R                [* -> 4.4.2]
## 
## - Lockfile written to "~/school_2025-2026/Biol430G/BIOL430G_mini-project/manuscript/docs/renv.lock".
\end{verbatim}

\begin{Shaded}
\begin{Highlighting}[]
\DocumentationTok{\#\#Figure 2}
\CommentTok{\#Look at participant individual change between active and sham for meanValenceRating}
\FunctionTok{ggplot}\NormalTok{(data, }\FunctionTok{aes}\NormalTok{(}\AttributeTok{x =}\NormalTok{ condition, }\AttributeTok{y =}\NormalTok{ meanValenceRating, }\AttributeTok{group =}\NormalTok{ participant)) }\SpecialCharTok{+}
  \FunctionTok{geom\_point}\NormalTok{() }\SpecialCharTok{+}
  \FunctionTok{geom\_line}\NormalTok{(}\AttributeTok{alpha =} \FloatTok{0.5}\NormalTok{) }\SpecialCharTok{+}
  \FunctionTok{facet\_wrap}\NormalTok{(}\SpecialCharTok{\textasciitilde{}}\NormalTok{stimuliType) }\SpecialCharTok{+}
  \FunctionTok{labs}\NormalTok{(}\AttributeTok{title =} \StringTok{"Individual Differences in Valence Ratings by Condition and Stimulus Type"}\NormalTok{) }\SpecialCharTok{+}
  \FunctionTok{theme\_minimal}\NormalTok{()}
\end{Highlighting}
\end{Shaded}

\pandocbounded{\includegraphics[keepaspectratio]{2026-03-12__manuscript_rmd_files/figure-latex/unnamed-chunk-3-1.pdf}}

\begin{Shaded}
\begin{Highlighting}[]
\DocumentationTok{\#\#Figure 3}
\CommentTok{\#Look at participant individual change between active and sham for meanArousalRating}
\FunctionTok{ggplot}\NormalTok{(data, }\FunctionTok{aes}\NormalTok{(}\AttributeTok{x =}\NormalTok{ condition, }\AttributeTok{y =}\NormalTok{ meanArousalRating, }\AttributeTok{group =}\NormalTok{ participant)) }\SpecialCharTok{+}
  \FunctionTok{geom\_point}\NormalTok{() }\SpecialCharTok{+}
  \FunctionTok{geom\_line}\NormalTok{(}\AttributeTok{alpha =} \FloatTok{0.5}\NormalTok{) }\SpecialCharTok{+}
  \FunctionTok{facet\_wrap}\NormalTok{(}\SpecialCharTok{\textasciitilde{}}\NormalTok{stimuliType) }\SpecialCharTok{+}
  \FunctionTok{labs}\NormalTok{(}\AttributeTok{title =} \StringTok{"Individual Differences in Arousal Ratings by Condition and Stimulus Type"}\NormalTok{) }\SpecialCharTok{+}
  \FunctionTok{theme\_minimal}\NormalTok{()}
\end{Highlighting}
\end{Shaded}

\pandocbounded{\includegraphics[keepaspectratio]{2026-03-12__manuscript_rmd_files/figure-latex/unnamed-chunk-4-1.pdf}}

\begin{Shaded}
\begin{Highlighting}[]
\DocumentationTok{\#\#Figure 4}
\FunctionTok{ggplot}\NormalTok{(data, }\FunctionTok{aes}\NormalTok{(}\AttributeTok{x =}\NormalTok{ GoldMSI\_ActiveEngagement\_63,}
                 \AttributeTok{y =}\NormalTok{ meanArousalRating,}
                 \AttributeTok{color =}\NormalTok{ stimuliType)) }\SpecialCharTok{+}
  \FunctionTok{geom\_point}\NormalTok{(}\AttributeTok{alpha =} \FloatTok{0.5}\NormalTok{) }\SpecialCharTok{+}
  \FunctionTok{geom\_smooth}\NormalTok{(}\AttributeTok{method =} \StringTok{"lm"}\NormalTok{, }\AttributeTok{se =} \ConstantTok{TRUE}\NormalTok{) }\SpecialCharTok{+}
  \FunctionTok{labs}\NormalTok{(}
    \AttributeTok{title =} \StringTok{"GoldMSI Active Engagement Predicting Arousal Ratings"}\NormalTok{,}
    \AttributeTok{x =} \StringTok{"GoldMSI Active Engagement"}\NormalTok{,}
    \AttributeTok{y =} \StringTok{"Mean Arousal Rating"}\NormalTok{,}
    \AttributeTok{color =} \StringTok{"Stimulus Type"}
\NormalTok{  ) }\SpecialCharTok{+}
  \FunctionTok{theme\_minimal}\NormalTok{()}
\end{Highlighting}
\end{Shaded}

\begin{verbatim}
## `geom_smooth()` using formula = 'y ~ x'
\end{verbatim}

\pandocbounded{\includegraphics[keepaspectratio]{2026-03-12__manuscript_rmd_files/figure-latex/unnamed-chunk-5-1.pdf}}

\begin{Shaded}
\begin{Highlighting}[]
\DocumentationTok{\#\#Figure 5}
\CommentTok{\#plot }
\FunctionTok{ggplot}\NormalTok{(data, }\FunctionTok{aes}\NormalTok{(}\AttributeTok{x =}\NormalTok{ GoldMSI\_Emotions\_42,}
                 \AttributeTok{y =}\NormalTok{ meanArousalRating,}
                 \AttributeTok{color =}\NormalTok{ stimuliType)) }\SpecialCharTok{+}
  \FunctionTok{geom\_point}\NormalTok{(}\AttributeTok{alpha =} \FloatTok{0.45}\NormalTok{, }\AttributeTok{size =} \DecValTok{2}\NormalTok{) }\SpecialCharTok{+}
  \FunctionTok{geom\_smooth}\NormalTok{(}\AttributeTok{method =} \StringTok{"lm"}\NormalTok{, }\AttributeTok{se =} \ConstantTok{TRUE}\NormalTok{) }\SpecialCharTok{+}
  \FunctionTok{labs}\NormalTok{(}
    \AttributeTok{title =} \StringTok{"Emotions and Arousal Ratings"}\NormalTok{,}
    \AttributeTok{x =} \StringTok{"GoldMSI Emotional Engagement"}\NormalTok{,}
    \AttributeTok{y =} \StringTok{"Mean Arousal Rating"}\NormalTok{,}
    \AttributeTok{color =} \StringTok{"Stimulus Type"}
\NormalTok{  ) }\SpecialCharTok{+}
  \FunctionTok{theme\_minimal}\NormalTok{(}\AttributeTok{base\_size =} \DecValTok{13}\NormalTok{)}
\end{Highlighting}
\end{Shaded}

\begin{verbatim}
## `geom_smooth()` using formula = 'y ~ x'
\end{verbatim}

\pandocbounded{\includegraphics[keepaspectratio]{2026-03-12__manuscript_rmd_files/figure-latex/unnamed-chunk-6-1.pdf}}

Note that the \texttt{echo\ =\ FALSE} parameter was added to the code
chunk to prevent printing of the R code that generated the plot.

\end{document}
